%% Formule de l'operateur : pos_sol_fixe_loa
%% Dans ce fichier, les �ventuels accents sont cod�s
%% en ISO-8859-15 (aka latin9)
%% La ligne suivante donne trois exemples 
%% e accent aigu : �, e accent grave : �, e-dans-l'o : �

Calcul des deux coordonn�es du soleil dans le rep�re �quatorial local :
\[
\begin{cases} \omega &=  15  \frac{\pi}{180} \left(\frac{Time}{3600} + \frac{Lon}{15} - \frac{ET}{60} - 12 \right)\\
\delta &= f\left(2\pi \frac{(JDA-1)}{365}\right)
\end{cases} 
\]
%�quation du temps
$\omega$: angle horaire, l'angle sur le plan �quatorial, 
entre la position du plan m�ridien contenant le soleil et le m�ridien local,
compt�  positivement vers l'Est.
%  	Angle horaire:
%L'une des coordonn�es horaires. Angle di�dre du cercle horaire de la direction et du m�ridien du lieu pris comme
%origine. L'angle horaire est compt� positivement dans le sens r�trograde. 

$\delta$ : d�clinaison, l'angle que forme la direction Terre-Soleil par rapport au 
plan de l'�quateur terrestre, dans le cercle m�ridien contenant le soleil.

$ET$ : l'��quation� du temps, la diff�rence (Temps solaire vrai) - (Temps solaire moyen).

Puis conversion g�om�trique en azimut ($\Phi_0$) et angle solaire z�nithal ($\Theta_0$):
\[
\begin{cases}
Elev &=  \arcsin \left( \sin(Lat) \sin(\delta) + \cos(Lat)\cos(\delta)\cos(\omega) \right)\\
 \Theta_0 &=  90 - Elev\\
\Phi_0 &\text{est tel que } \begin{cases}\sin(\Phi_O)&= \cos(\delta)*\sin(\omega)/\cos(Elev) \\
\cos(\Phi_O)&= \frac{-\cos(Lat)\sin(\delta)+\sin(Lat)\cos(\delta)\cos(\omega)}{\cos(Elev)}
\end{cases} 
\end{cases} 
\]
